
\documentclass[11pt,a4paper]{article}

% =========================================================
% PACKAGES
% =========================================================

\usepackage[T1]{fontenc}
\usepackage[utf8]{inputenc}
\usepackage[english]{babel}
\usepackage{lmodern,microtype,amsmath}

\usepackage[
    a4paper,
    left=24mm,
    right=24mm,
    top=24mm,
    bottom=24mm,
    headheight=14pt
]{geometry}

\usepackage[table]{xcolor}
\usepackage{pagecolor}
\usepackage{graphicx,float,caption,booktabs,tabularx,array,longtable}
\usepackage{fancyhdr,lastpage,titlesec,enumitem,listings,hyperref}
\usepackage{etoolbox}
\usepackage{needspace}

% =========================================================
% DESIGN AND COLOURS
% =========================================================

\definecolor{paperpink}{HTML}{FFF9FB}
\definecolor{ink}{HTML}{222222}
\definecolor{muted}{HTML}{555555}
\definecolor{rulegray}{HTML}{D8D5D7}
\definecolor{codebg}{HTML}{FAF8F9}

\pagecolor{paperpink}
\color{ink}

\hypersetup{
    colorlinks=true,
    allcolors=ink,
    pdftitle={White Wine Quality Prediction},
    pdfauthor={Gabriela Rehakova}
}

\setlength{\parindent}{0pt}
\setlength{\parskip}{5pt}
\setlength{\emergencystretch}{2em}

\renewcommand{\arraystretch}{1.16}

\setlist[itemize]{
    topsep=3pt,
    itemsep=2pt,
    leftmargin=1.5em
}

\setlist[enumerate]{
    topsep=3pt,
    itemsep=2pt,
    leftmargin=1.5em
}

% =========================================================
% SECTION FORMATTING
% =========================================================

\titleformat{\section}
    {\Large\bfseries\color{ink}}
    {\thesection}{0.7em}{}
    [\vspace{2pt}\color{rulegray}\titlerule]

\titleformat{\subsection}
    {\large\bfseries\color{ink}}
    {\thesubsection}{0.65em}{}

\titlespacing*{\section}{0pt}{17pt}{9pt}
\titlespacing*{\subsection}{0pt}{11pt}{4pt}

\pretocmd{\section}{\Needspace{6\baselineskip}}{}{}
\pretocmd{\subsection}{\Needspace{4\baselineskip}}{}{}

% =========================================================
% HEADER AND FOOTER
% =========================================================

\pagestyle{fancy}
\fancyhf{}

\fancyhead[L]{\footnotesize White Wine Quality Prediction}
\fancyhead[R]{\footnotesize SEPL Technical Report}

\fancyfoot[L]{\footnotesize Gabriela Rehakova BSc}
\fancyfoot[R]{\footnotesize\thepage}

\renewcommand{\headrulewidth}{0.35pt}
\renewcommand{\footrulewidth}{0pt}

% =========================================================
% FIGURES AND CODE LISTINGS
% =========================================================

\captionsetup{
    font=small,
    labelfont=bf,
    textfont=normalfont,
    skip=5pt
}

\lstset{
    basicstyle=\ttfamily\footnotesize,
    backgroundcolor=\color{codebg},
    frame=single,
    rulecolor=\color{rulegray},
    framerule=.4pt,
    breaklines=true,
    showstringspaces=false,
    columns=fullflexible,
    aboveskip=6pt,
    belowskip=6pt
}

\graphicspath{{../screenshots/}}

% Required screenshot.
% Missing files are displayed as placeholders.

\newcommand{\shot}[3]{%
    \begin{figure}[H]
        \centering
        \IfFileExists{../screenshots/#1}{%
            \includegraphics[
                width=.94\linewidth,
                height=.55\textheight,
                keepaspectratio
            ]{../screenshots/#1}
        }{%
            \fbox{%
                \parbox[c][25mm][c]{.86\linewidth}{%
                    \centering\small
                    Screenshot:
                    \texttt{\detokenize{#1}}\\
                    Upload this file to docs/screenshots/
                }
            }
        }
        \caption{#2}
        \label{#3}
    \end{figure}
}

% Optional screenshot.
% Included only when the corresponding file exists.

\newcommand{\optionalshot}[3]{%
    \IfFileExists{../screenshots/#1}{%
        \shot{#1}{#2}{#3}
    }{}
}

\newcommand{\codepath}[1]{\texttt{\detokenize{#1}}}

% =========================================================
% DOCUMENT
% =========================================================

\begin{document}

% =========================================================
% TITLE PAGE
% =========================================================

\begin{titlepage}
\thispagestyle{empty}

\vspace*{8mm}

{\small\bfseries FH TECHNIKUM WIEN}\par

\vspace{3mm}

{\small Master's in AI \& Data Science}\par

\vspace{15mm}

\rule{40mm}{1pt}\par

\vspace{10mm}

{\fontsize{31}{36}\selectfont\bfseries
White Wine\\[3mm]
Quality Prediction\par}

\vspace{8mm}

{\large A Containerized Machine Learning System\par}

\vspace{4mm}

{\normalsize
Asynchronous Inference, Automated Testing,
Continuous Integration and Application Monitoring\par}

\vspace{16mm}

\begin{tabularx}{\linewidth}{@{}XXX@{}}
    \textbf{75.16\%} &
    \textbf{23 / 23} &
    \textbf{3}\\
    \small Model accuracy &
    \small Tests passed &
    \small Docker services
\end{tabularx}

\vspace{8mm}

\begin{tabularx}{\linewidth}{@{}XXX@{}}
    \textbf{2 / 2} &
    \textbf{100\%} &
    \textbf{1m 40s}\\
    \small CI jobs passed &
    \small Test success rate &
    \small Latest CI duration
\end{tabularx}

\vspace{12mm}

\rule{\linewidth}{.5pt}\par

\vspace{6mm}

\renewcommand{\arraystretch}{1.4}

\begin{tabularx}{\linewidth}{@{}p{35mm}X@{}}
    \textbf{Author} & Gabriela Rehakova BSc\\
    \textbf{Course} & Software Engineering Practices in ML\\
    \textbf{Programme} & Master's in AI \& Data Science\\
    \textbf{Academic year} & 2026/2027\\
    \textbf{Date} & 24 November 2026\\
    \textbf{Repository} &
    \url{https://github.com/neskrotna/sepl-ml-system}
\end{tabularx}

\vfill

{\small\color{muted}
Technical project report
\hfill FH Technikum Wien}\par

\end{titlepage}

\pagenumbering{arabic}

% =========================================================
% EXECUTIVE SUMMARY
% =========================================================

\section*{Executive Summary}
\addcontentsline{toc}{section}{Executive Summary}

This project documents the development of a containerized
machine learning system for binary white wine quality
prediction.

A Random Forest classifier was trained using physicochemical
features from the UCI White Wine Quality dataset. After
removing 937 duplicate records, 3,961 observations remained.
The trained model achieved a test accuracy of
\textbf{75.16\%}.

The prediction system combines a FastAPI REST interface,
Pydantic input validation, Redis-based asynchronous task
processing and a Python background worker.

Docker Compose orchestrates the API, Redis and worker
services. The infrastructure includes environment-based
configuration, read-only model bind mounts, persistent
Redis storage, service networking and Redis health checks.

Application logging records startup events, Redis
connectivity, prediction requests, task processing,
prediction results and inference execution time.

The system passed all \textbf{23 automated tests} locally,
including model unit tests, API tests and end-to-end
integration tests. Ruff reported no remaining lint issues.

Continuous integration was implemented using GitHub Actions.
An initial account-level billing restriction was resolved
by GitHub Support. Subsequent workflow executions completed
successfully on GitHub-hosted runners.

The latest verified workflow was automatically triggered
by a push to the main branch. Both CI jobs passed, with a
total execution duration of \textbf{1 minute and 40 seconds}.

The implementation demonstrates a reproducible ML
engineering workflow from model development to
containerized inference, automated testing and continuous
integration.

The HPC simulation and final performance reflection remain
outstanding at the current project stage.

\begin{table}[H]
\centering
\begin{tabularx}{\linewidth}{@{}lXlX@{}}
\toprule
\textbf{Model} & Random Forest &
\textbf{API} & FastAPI\\

\textbf{Dataset} & UCI White Wine &
\textbf{Queue} & Redis\\

\textbf{Infrastructure} & Docker Compose &
\textbf{Testing} & pytest, Ruff\\

\textbf{CI} & GitHub Actions &
\textbf{Logging} & Python logging\\

\textbf{Accuracy} & 75.16\% &
\textbf{Tests} & 23/23 passed\\
\bottomrule
\end{tabularx}
\caption{Project summary.}
\end{table}

\clearpage
\tableofcontents
\clearpage

% =========================================================
% 1. INTRODUCTION
% =========================================================

\section{Introduction}

The objective of this individual assessment was to engineer
a functioning machine learning inference application rather
than deliver a standalone predictive model.

The project combines machine learning development with
software engineering practices covering application
architecture, API design, asynchronous processing,
containerization, testing, continuous integration and
application monitoring.

The selected use case is binary white wine quality
prediction based on eleven numerical physicochemical
measurements.

The system accepts prediction requests through a REST API,
validates the supplied features, queues prediction jobs
and performs inference using a background worker.

The implementation separates model training from runtime
inference. A trained and serialized model artifact is
reused by the application without requiring retraining
for each prediction.

\subsection{Project Objectives}

The main project objectives are:

\begin{itemize}
    \item Develop and evaluate a machine learning model.
    \item Persist the trained model for reusable inference.
    \item Implement a validated REST API.
    \item Separate request handling from inference.
    \item Implement Redis-based asynchronous job processing.
    \item Containerize the application services.
    \item Externalize environment configuration.
    \item Manage the model through read-only bind mounts.
    \item Implement automated tests and static code analysis.
    \item Configure continuous integration.
    \item Implement application logging and health monitoring.
    \item Document implementation and verification results.
\end{itemize}

\subsection{Technology Stack}

\begin{table}[H]
\centering
\begin{tabularx}{\linewidth}{@{}lX@{}}
\toprule
\textbf{Area} & \textbf{Technologies}\\
\midrule

ML development &
Python, pandas, NumPy, scikit-learn\\

Model persistence &
Joblib\\

API and validation &
FastAPI, Pydantic, Uvicorn\\

Queue and worker &
Redis, Python\\

Infrastructure &
Docker, Docker Compose\\

Configuration &
Environment variables, .env files\\

Testing &
pytest, HTTPX, FastAPI TestClient\\

Static analysis &
Ruff\\

Version control &
Git, GitHub\\

Continuous integration &
GitHub Actions\\

Logging and monitoring &
Python logging, Docker logs, health checks\\

Documentation &
LaTeX, Markdown, Overleaf\\

\bottomrule
\end{tabularx}
\caption{Project technology stack.}
\end{table}

% =========================================================
% 2. DATASET AND PREPROCESSING
% =========================================================

\section{Dataset and Preprocessing}

\subsection{Dataset Description}

The project uses the UCI White Wine Quality dataset
\cite{wine,uci}.

The original dataset contains 4,898 observations,
eleven numerical physicochemical attributes and an
original wine quality score.

The dataset is stored in:

\begin{lstlisting}
data/winequality-white.csv
\end{lstlisting}

\begin{table}[H]
\centering
\begin{tabularx}{\linewidth}{@{}lX@{}}
\toprule
\textbf{Feature} & \textbf{Description}\\
\midrule
Fixed acidity & Non-volatile acidity\\
Volatile acidity & Volatile acid concentration\\
Citric acid & Citric acid concentration\\
Residual sugar & Remaining sugar after fermentation\\
Chlorides & Salt concentration\\
Free sulfur dioxide & Free sulfur dioxide concentration\\
Total sulfur dioxide & Total sulfur dioxide concentration\\
Density & Wine density\\
pH & Acidity level\\
Sulphates & Sulphate concentration\\
Alcohol & Alcohol content\\
\bottomrule
\end{tabularx}
\caption{Input features used for wine quality prediction.}
\end{table}

\subsection{Exploratory Data Analysis}

The dataset was explored before model training to
understand its structure, feature distributions and
data quality.

\shot{01_dataset_exploration.png}
{Initial exploration of the white wine dataset.}
{fig:data1}

\shot{02_dataset_exploration.png}
{Additional exploratory analysis of the dataset.}
{fig:data2}

\subsection{Data Cleaning}

Duplicate removal reduced the dataset from 4,898 to
3,961 observations.

A total of 937 duplicate records were removed before
the training and test datasets were created.

Removing duplicates reduces the risk of identical
observations appearing in both subsets and influencing
the evaluation results.

\subsection{Binary Target Definition}

The original quality score was converted into a binary
classification target.

\[
y =
\begin{cases}
1, & \text{if quality} \geq 6\\
0, & \text{if quality} < 6
\end{cases}
\]

The resulting classes are:

\begin{table}[H]
\centering
\begin{tabularx}{\linewidth}{@{}clX@{}}
\toprule
\textbf{Class ID} & \textbf{Label} & \textbf{Definition}\\
\midrule
0 & \texttt{lower\_quality} & Quality score below 6\\
1 & \texttt{good\_quality} & Quality score of at least 6\\
\bottomrule
\end{tabularx}
\caption{Binary classification target.}
\end{table}

\subsection{Training and Test Split}

The cleaned dataset was divided into training and
test subsets.

\begin{table}[H]
\centering
\begin{tabular}{@{}lr@{}}
\toprule
\textbf{Stage} & \textbf{Observations}\\
\midrule
Original dataset & 4,898\\
Removed duplicates & 937\\
Cleaned dataset & 3,961\\
Training dataset & 3,168\\
Test dataset & 793\\
\bottomrule
\end{tabular}
\caption{Data preparation summary.}
\end{table}

The training dataset was used to fit the classifier,
while the test dataset was reserved for evaluating
predictive performance.

% =========================================================
% 3. MODEL TRAINING AND EVALUATION
% =========================================================

\section{Model Training and Evaluation}

\subsection{Random Forest Classifier}

A Random Forest classifier was selected because of its
suitability for numerical tabular data and its ability
to represent nonlinear relationships between features.

Random Forest combines multiple decision trees to
produce a final classification and supports probability
estimates for individual predictions.

The implementation uses scikit-learn \cite{sklearn}.

The training script is located at:

\begin{lstlisting}
ml_model/train.py
\end{lstlisting}

The script performs dataset loading, preprocessing,
target conversion, training, evaluation and model
serialization.

The model was trained using:

\begin{lstlisting}
python ml_model/train.py
\end{lstlisting}

\shot{03_model_training.png}
{Training the Random Forest classifier.}
{fig:train}

\subsection{Model Performance}

The trained classifier achieved a test accuracy of
\textbf{75.16\%}.

Accuracy is defined as:

\[
Accuracy =
\frac{TP+TN}{TP+TN+FP+FN}
\]

The confusion matrix, using class order 0 followed
by class 1, was:

\[
\begin{bmatrix}
162 & 108\\
89 & 434
\end{bmatrix}
\]

\begin{table}[H]
\centering
\begin{tabular}{@{}lrr@{}}
\toprule
& \textbf{Predicted 0} & \textbf{Predicted 1}\\
\midrule
Actual 0 & 162 & 108\\
Actual 1 & 89 & 434\\
\bottomrule
\end{tabular}
\caption{Random Forest confusion matrix.}
\end{table}

The macro-averaged F1-score was approximately 0.72.

Precision, recall and F1-score complement accuracy by
providing information about classification performance
for individual classes.

\[
Precision = \frac{TP}{TP+FP}
\]

\[
Recall = \frac{TP}{TP+FN}
\]

\[
F_1 =
2 \cdot
\frac{Precision \cdot Recall}
{Precision+Recall}
\]

\shot{04_model_evaluation.png}
{Model evaluation results.}
{fig:evaluation}

\subsection{Model Serialization}

The trained model was serialized using joblib and saved
as:

\begin{lstlisting}
ml_model/model.pkl
\end{lstlisting}

The serialized artifact has a size of approximately
7.47 MB.

Persisting the trained model separates model training
from application inference and avoids retraining during
application startup or prediction processing.

\shot{05_saved_model.png}
{Persisted Random Forest model artifact.}
{fig:artifact}

\subsection{Local Inference}

The model module is responsible for loading the trained
artifact, validating input features, preparing numerical
input and performing inference.

A local sample produced the following prediction values:

\begin{lstlisting}
{
    "prediction": "good_quality",
    "class_id": 1,
    "confidence": 0.75
}
\end{lstlisting}

The confidence value represents the model's estimated
probability for the predicted class. It should not be
interpreted as a guarantee of prediction correctness.

\shot{06_local_prediction.png}
{Successful local model inference.}
{fig:local}

% =========================================================
% 4. SYSTEM ARCHITECTURE
% =========================================================

\section{System Architecture}

\subsection{Architectural Overview}

The application consists of three cooperating services:

\begin{itemize}
    \item FastAPI REST API.
    \item Redis job queue and result store.
    \item Python background inference worker.
\end{itemize}

The API is responsible for receiving and validating
requests. Redis stores pending prediction jobs and task
results. The worker consumes queued jobs and performs
machine learning inference.

\begin{center}
\fbox{
\begin{minipage}{.86\linewidth}
\centering
\vspace{5pt}

\textbf{Client / Swagger UI}\\[4pt]
$\downarrow$\\[4pt]

\textbf{FastAPI REST API}\\
\small Request validation and task creation\\[4pt]
$\downarrow$\\[4pt]

\textbf{Redis Job Queue}\\
\small Pending prediction tasks\\[4pt]
$\downarrow$\\[4pt]

\textbf{Background Worker}\\
\small Input preprocessing and model inference\\[4pt]
$\downarrow$\\[4pt]

\textbf{Random Forest Model}\\[4pt]
$\downarrow$\\[4pt]

\textbf{Redis Result Store}\\[4pt]
$\downarrow$\\[4pt]

\textbf{FastAPI Result Endpoint}\\[4pt]
$\downarrow$\\[4pt]

\textbf{Client}

\vspace{5pt}
\end{minipage}
}
\end{center}

\subsection{Asynchronous Prediction Workflow}

The prediction pipeline follows these steps:

\begin{enumerate}
    \item The client submits a prediction request.
    \item FastAPI validates the request using Pydantic.
    \item The API generates a unique task identifier.
    \item The prediction job is added to Redis.
    \item The API returns HTTP 202 Accepted.
    \item The worker retrieves the queued job.
    \item The worker preprocesses the input features.
    \item The trained model performs inference.
    \item The prediction result is stored in Redis.
    \item The client retrieves the result using the task ID.
\end{enumerate}

The architecture decouples HTTP request handling from
machine learning inference.

For the selected lightweight classifier, the primary
benefit is separation of responsibilities and
asynchronous task management rather than a demonstrated
improvement in inference speed.

% =========================================================
% 5. REST API IMPLEMENTATION
% =========================================================

\section{REST API Implementation}

\subsection{FastAPI and Pydantic}

The REST API was implemented using FastAPI
\cite{fastapi}.

Pydantic schemas define the expected prediction features
and support automatic request validation.

The relevant implementation files are:

\begin{lstlisting}
api/main.py
api/schemas.py
\end{lstlisting}

Interactive API documentation is available through
Swagger UI:

\begin{center}
\url{http://localhost:8000/docs}
\end{center}

\begin{table}[H]
\centering
\begin{tabularx}{\linewidth}{@{}llX@{}}
\toprule
\textbf{Method} & \textbf{Endpoint} & \textbf{Purpose}\\
\midrule
GET & \texttt{/} & Application information\\
GET & \texttt{/health} & Redis connectivity check\\
POST & \texttt{/predict} & Submit a prediction job\\
GET & \texttt{/result/\{task\_id\}} & Retrieve task status or result\\
\bottomrule
\end{tabularx}
\caption{Implemented REST API endpoints.}
\end{table}

\shot{07_swagger_ui.png}
{Swagger UI and automatically generated API documentation.}
{fig:swagger}

\subsection{Prediction Request}

The client submits eleven physicochemical features
through the prediction endpoint.

An example request is:

\begin{lstlisting}
{
    "fixed_acidity": 7.0,
    "volatile_acidity": 0.27,
    "citric_acid": 0.36,
    "residual_sugar": 20.7,
    "chlorides": 0.045,
    "free_sulfur_dioxide": 45.0,
    "total_sulfur_dioxide": 170.0,
    "density": 1.001,
    "pH": 3.0,
    "sulphates": 0.45,
    "alcohol": 8.8
}
\end{lstlisting}

A successfully accepted request returns HTTP 202 and
a generated task identifier.

\shot{08_predict_request.png}
{Prediction request submitted through the REST API.}
{fig:predict}

\shot{09_queued_result.png}
{Queued prediction status returned by the API.}
{fig:queued}

\subsection{Prediction Result Retrieval}

The result endpoint retrieves the current task state
using its unique identifier.

While processing is pending, the task remains queued.
After inference finishes, the response contains the
completed prediction information.

\shot{14_new_prediction_request.png}
{Additional prediction request submitted through Swagger UI.}
{fig:second}

\shot{15_completed_api_prediction.png}
{Completed prediction retrieved through the result endpoint.}
{fig:completed}

\subsection{Input Validation and Error Handling}

Pydantic validates request structure and expected
feature types.

Additional validation in the model layer checks missing
features, unexpected features, nonnumeric values,
NaN values and infinite values.

The API implements appropriate HTTP responses for
successful requests and error conditions.

\begin{table}[H]
\centering
\begin{tabularx}{\linewidth}{@{}lX@{}}
\toprule
\textbf{HTTP Status} & \textbf{Meaning}\\
\midrule
200 OK & Successful request\\
202 Accepted & Prediction job accepted\\
404 Not Found & Unknown task identifier\\
422 Unprocessable Entity & Invalid prediction input\\
503 Service Unavailable & Redis connection unavailable\\
\bottomrule
\end{tabularx}
\caption{HTTP response codes.}
\end{table}

\shot{16_api_validation_error.png}
{Validation error for invalid prediction input.}
{fig:validation}

\shot{17_unknown_task_error.png}
{Unknown task identifier returning HTTP 404.}
{fig:unknown}

% =========================================================
% 6. REDIS AND WORKER PROCESSING
% =========================================================

\section{Redis and Worker Processing}

\subsection{Redis Job Queue}

Redis is used as a lightweight job queue and result
store \cite{redis}.

Prediction jobs are placed into the following Redis list:

\begin{lstlisting}
wine_prediction_jobs
\end{lstlisting}

Task results are stored using keys with the following
structure:

\begin{lstlisting}
result:<task_id>
\end{lstlisting}

The use of unique task identifiers supports independent
tracking of submitted prediction jobs.

\shot{10_redis_queue.png}
{Prediction job stored in the Redis queue.}
{fig:redisqueue}

\subsection{Background Worker}

The background worker is implemented in:

\begin{lstlisting}
worker/worker.py
\end{lstlisting}

The worker loads the trained Random Forest model,
establishes a Redis connection and waits for prediction
jobs.

After receiving a job, the worker preprocesses the
features, performs inference and stores the result
in Redis.

\shot{11_worker_processing.png}
{Background worker processing a queued prediction.}
{fig:worker}

\subsection{Queue Consumption and Result Storage}

Successful job consumption was verified by inspecting
the Redis queue after worker processing.

\shot{12_redis_empty_queue.png}
{Redis queue after successful job consumption.}
{fig:empty}

The completed prediction was subsequently verified
in the Redis result store.

\shot{13_redis_completed_result.png}
{Completed prediction result stored in Redis.}
{fig:stored}

% =========================================================
% 7. DOCKER CONTAINERIZATION
% =========================================================

\section{Docker Containerization}

\subsection{Containerized Infrastructure}

Docker was used to package the application and its
runtime dependencies into reproducible containers
\cite{docker}.

Docker Compose coordinates the API, Redis and worker
services.

The infrastructure is defined in:

\begin{lstlisting}
Dockerfile
compose.yaml
.dockerignore
.env.example
\end{lstlisting}

The Dockerfile packages the Python application and
required dependencies.

The Compose configuration defines service commands,
environment variables, networking, volumes, restart
policies and startup dependencies.

\begin{table}[H]
\centering
\begin{tabularx}{\linewidth}{@{}lX@{}}
\toprule
\textbf{Service} & \textbf{Responsibility}\\
\midrule
API & HTTP request handling and validation\\
Redis & Job queue and result storage\\
Worker & Background machine learning inference\\
\bottomrule
\end{tabularx}
\caption{Docker Compose services.}
\end{table}

\subsection{Initial Container Deployment}

The initial application deployment was performed using:

\begin{lstlisting}
docker compose up -d --build
docker compose ps
docker compose logs worker
\end{lstlisting}

\shot{18_docker_compose_build.png}
{Initial Docker Compose image build.}
{fig:build}

\shot{19_docker_compose_services.png}
{Initial running Docker Compose services.}
{fig:services}

\shot{20_docker_worker_logs.png}
{Background worker container logs.}
{fig:logs}

\subsection{Initial End-to-End Verification}

The containerized application was verified by submitting
prediction requests through the REST API.

Redis accepted queued jobs, the background worker
performed inference and the result endpoint returned
completed predictions.

\shot{21_docker_prediction_request.png}
{Prediction request submitted to the Dockerized API.}
{fig:drequest}

\shot{22_docker_completed_prediction.png}
{Completed prediction returned by the Dockerized API.}
{fig:dcompleted}

\shot{23_docker_worker_inference.png}
{Successful worker inference inside Docker.}
{fig:dinference}

\shot{24_docker_redis_verification.png}
{Redis verification inside the Docker environment.}
{fig:dredis}

\subsection{Environment-Based Configuration}

The infrastructure was subsequently extended to
externalize runtime configuration.

The following environment variables define the Redis
connection settings:

\begin{lstlisting}
REDIS_HOST=redis
REDIS_PORT=6379
\end{lstlisting}

A version-controlled configuration template was created:

\begin{lstlisting}
.env.example
\end{lstlisting}

The local runtime configuration is stored in:

\begin{lstlisting}
.env
\end{lstlisting}

The actual environment file is excluded from Git,
while the example file is included in the repository.

On Windows, the configuration is created using:

\begin{lstlisting}
Copy-Item .env.example .env
\end{lstlisting}

On Linux and macOS:

\begin{lstlisting}
cp .env.example .env
\end{lstlisting}

Docker Compose passes the configuration to the API
and worker containers.

Externalizing configuration separates environment-specific
settings from application source code and supports
reproducible deployments.

\subsection{Read-Only Model Bind Mounts}

The trained Random Forest artifact was configured as a
read-only bind mount in both the API and worker
containers.

\begin{lstlisting}
volumes:
  - ./ml_model/model.pkl:/app/ml_model/model.pkl:ro
\end{lstlisting}

The host model artifact is mounted at:

\begin{lstlisting}
/app/ml_model/model.pkl
\end{lstlisting}

The read-only option prevents the containers from
modifying the mounted model file.

This configuration supports independent runtime
management of the trained artifact. Replacing the
host model file does not require rebuilding the
application image, although the worker must reload
the model before using the replacement.

The Dockerfile also packages the model directory.
The runtime bind mount overrides the corresponding
model file inside the container.

\subsection{Service Networking}

Docker Compose creates a shared default network for
the application services.

The API and worker communicate with Redis using its
service hostname:

\begin{lstlisting}
redis:6379
\end{lstlisting}

The Redis service does not publish its port to the
host machine.

The API port is published through the loopback
interface:

\begin{lstlisting}
127.0.0.1:8000:8000
\end{lstlisting}

This configuration makes the API accessible locally
while avoiding direct exposure through all host
network interfaces.

\subsection{Redis Data Persistence}

Redis uses a named Docker volume for persistent
storage.

\begin{lstlisting}
volumes:
  - redis_data:/data
\end{lstlisting}

Append-only persistence is enabled:

\begin{lstlisting}
command: ["redis-server", "--appendonly", "yes"]
\end{lstlisting}

The named volume preserves Redis data across normal
container shutdown and recreation.

This configuration does not provide a complete backup
or disaster recovery mechanism.

\subsection{Redis Health Check}

A Redis health check was configured using the
Redis CLI.

\begin{lstlisting}
healthcheck:
  test: ["CMD", "redis-cli", "ping"]
  interval: 5s
  timeout: 3s
  retries: 10
\end{lstlisting}

The API and worker depend on Redis becoming healthy
before starting.

\begin{lstlisting}
depends_on:
  redis:
    condition: service_healthy
\end{lstlisting}

This reduces startup failures caused by attempting to
connect to Redis before it is ready.

\subsection{Restart Policies}

The application services use the following restart
policy:

\begin{lstlisting}
restart: unless-stopped
\end{lstlisting}

This supports automatic container restarts following
unexpected termination, unless the container was
explicitly stopped.

\subsection{Updated Compose Configuration Verification}

The resolved Compose configuration was inspected using:

\begin{lstlisting}
docker compose config
\end{lstlisting}

The command verified the service definitions,
environment variables, model bind mounts, networking,
Redis persistence and health-check configuration.

\shot{37_docker_compose_configuration1.png}
{Resolved Docker Compose configuration, part one.}
{fig:composeconfig1}

\shot{37_docker_compose_configuration2.png}
{Resolved Docker Compose configuration, part two.}
{fig:composeconfig2}

\shot{37_docker_compose_configuration3.png}
{Resolved Docker Compose configuration, part three.}
{fig:composeconfig3}

\subsection{Updated Container Deployment}

The updated infrastructure was started and inspected
using:

\begin{lstlisting}
docker compose up -d --build
docker compose ps
\end{lstlisting}

All three services were running successfully.

The Redis container was marked as healthy, and the API
was accessible through its configured local port.

\shot{38_updated_docker_services.png}
{Running services after the infrastructure configuration update.}
{fig:updatedservices}

\subsection{Runtime Bind Mount Verification}

The model bind mounts were inspected directly in the
running API and worker containers.

The following commands were executed in PowerShell:

\begin{lstlisting}
docker inspect (docker compose ps -q api) --format '{{json .Mounts}}'

docker inspect (docker compose ps -q worker) --format '{{json .Mounts}}'
\end{lstlisting}

The inspection confirmed the following mount
properties:

\begin{table}[H]
\centering
\begin{tabularx}{\linewidth}{@{}lX@{}}
\toprule
\textbf{Property} & \textbf{Observed Value}\\
\midrule
Type & Bind mount\\
Destination & \codepath{/app/ml_model/model.pkl}\\
Mode & Read-only (ro)\\
RW & False\\
Services & API and worker\\
\bottomrule
\end{tabularx}
\caption{Verified model bind mount properties.}
\end{table}

\shot{39_model_bind_mount_verification.png}
{Read-only model bind mounts verified in both containers.}
{fig:bindmount}

\subsection{Runtime Environment Verification}

The environment variables were inspected inside both
application containers:

\begin{lstlisting}
docker compose exec api printenv REDIS_HOST REDIS_PORT

docker compose exec worker printenv REDIS_HOST REDIS_PORT
\end{lstlisting}

Both containers returned:

\begin{lstlisting}
redis
6379
\end{lstlisting}

This confirmed that the expected runtime configuration
was available inside the running services.

\shot{40_environment_variables_verification.png}
{Redis environment variables verified in both containers.}
{fig:environment}

\subsection{Regression Verification of Prediction Processing}

After introducing the new infrastructure configuration,
the end-to-end prediction workflow was executed again.

The API accepted a prediction request, the worker
loaded the mounted model and performed inference,
and the completed result was retrieved successfully.

\shot{41_prediction_after_configuration_update.png}
{Successful prediction after the Docker configuration update.}
{fig:updatedprediction}

\subsection{Containerization Verification Summary}

\begin{table}[H]
\centering
\begin{tabularx}{\linewidth}{@{}Xl@{}}
\toprule
\textbf{Verification Item} & \textbf{Result}\\
\midrule
Docker Compose configuration & Verified\\
API container startup & Successful\\
Worker container startup & Successful\\
Redis container startup & Successful\\
Redis health check & Healthy\\
Model bind mounts & Verified\\
Read-only model access & Verified\\
Environment variables & Verified\\
Service networking & Operational\\
Redis named volume configuration & Verified\\
End-to-end prediction & Successful\\
Automated integration tests & Passed\\
GitHub Actions integration job & Passed\\
\bottomrule
\end{tabularx}
\caption{Containerization verification summary.}
\end{table}

% =========================================================
% 8. AUTOMATED TESTING
% =========================================================

\section{Automated Testing}

\subsection{Testing Strategy}

The project uses pytest \cite{pytest} to verify
machine learning functionality, API behaviour and
integration between application services.

The testing strategy consists of three levels:

\begin{itemize}
    \item Model unit tests.
    \item API tests using mocked Redis interactions.
    \item End-to-end integration tests.
\end{itemize}

The complete test suite contains 23 automated tests.

\subsection{Model Unit Tests}

Ten tests are implemented in:

\begin{lstlisting}
tests/test_model.py
\end{lstlisting}

The tests cover:

\begin{itemize}
    \item Successful model loading.
    \item Valid input preprocessing.
    \item Missing input features.
    \item Unexpected input features.
    \item Nonnumeric input values.
    \item NaN input values.
    \item Infinite input values.
    \item Prediction response structure.
    \item Prediction confidence.
    \item Prediction for a known wine sample.
\end{itemize}

The tests were executed using:

\begin{lstlisting}
python -m pytest tests/test_model.py -v
\end{lstlisting}

\shot{25_pytest_unit_tests.png}
{Ten successfully executed model unit tests.}
{fig:unit}

\subsection{API Tests}

Eleven tests are implemented in:

\begin{lstlisting}
tests/test_api.py
\end{lstlisting}

The tests use FastAPI TestClient and mocked Redis
operations.

The tested scenarios include:

\begin{itemize}
    \item Root endpoint.
    \item Successful health check.
    \item Redis unavailability during health checks.
    \item Valid prediction requests.
    \item Missing prediction features.
    \item Invalid prediction feature types.
    \item Redis unavailability during prediction submission.
    \item Queued prediction results.
    \item Completed prediction results.
    \item Unknown task identifiers.
    \item Redis unavailability during result retrieval.
\end{itemize}

The tests were executed using:

\begin{lstlisting}
python -m pytest tests/test_api.py -v
\end{lstlisting}

\shot{26_pytest_api_tests.png}
{Eleven successfully executed API tests.}
{fig:apiunit}

\subsection{Integration Tests}

Two integration tests are implemented in:

\begin{lstlisting}
tests/test_integration.py
\end{lstlisting}

Unlike the isolated tests, integration tests communicate
with the running Docker application.

The first test verifies the complete prediction pipeline
from request submission to result retrieval.

The second test verifies HTTP 404 behaviour for an
unknown task identifier.

\begin{lstlisting}
python -m pytest tests/test_integration.py -v
\end{lstlisting}

\shot{27_pytest_integration_tests.png}
{Two successfully executed Docker integration tests.}
{fig:integration}

\subsection{Complete Test Suite}

The complete test suite was executed using:

\begin{lstlisting}
python -m pytest tests/ -v
\end{lstlisting}

All 23 tests passed.

\shot{28_pytest_full_suite.png}
{Initial complete test suite with 23 passing tests.}
{fig:full}

\subsection{Isolated Test Execution}

The unit and API tests were also executed independently
of the Docker integration tests.

\begin{lstlisting}
python -m pytest tests/ -v -m "not integration"
\end{lstlisting}

The execution produced 21 passing tests and two
deselected integration tests.

\shot{29_pytest_without_integration.png}
{Isolated unit and API test execution.}
{fig:nointegration}

\subsection{Local Environment Migration}

During project development, the application was
transferred to a new Windows development computer.

The previously copied virtual environment contained
references to the Python interpreter installed on
the original computer.

Consequently, local Python commands initially failed
because the referenced interpreter was no longer
available.

The virtual environment was removed and recreated
using the Python installation available on the new
computer.

\begin{lstlisting}
Remove-Item -Recurse -Force .venv

python -m venv .venv

.\.venv\Scripts\Activate.ps1

python -m pip install --upgrade pip

python -m pip install -r requirements.txt
\end{lstlisting}

The recreated environment used Python 3.13.14.

The Docker and GitHub Actions environments continued
to use Python 3.12.

This migration demonstrated the importance of
recreating virtual environments from dependency
definitions rather than transferring existing
virtual environment directories between computers.

\subsection{Model Loading Diagnostics}

During the initial regression test execution on the
new computer, the model-loading test appeared to
stall.

The affected test was isolated and executed with
additional diagnostic output:

\begin{lstlisting}
python -m pytest tests/test_model.py::test_model_loads_successfully -v -s -o faulthandler_timeout=30
\end{lstlisting}

The isolated test subsequently passed in 8.60 seconds.

The exact cause of the earlier delay was not
established.

No model retraining or application code changes were
required to resolve the observed behaviour.

\subsection{Regression Testing After Configuration Changes}

Following the introduction of environment-based
configuration and read-only model bind mounts,
the complete test suite was executed again.

\begin{lstlisting}
python -m pytest tests/ -v -o faulthandler_timeout=30
\end{lstlisting}

All 23 automated tests passed successfully.

The result confirmed that the updated infrastructure
did not introduce regressions detectable by the
existing test suite.

\shot{42_tests_after_configuration_update.png}
{All 23 tests passing after infrastructure changes and environment migration.}
{fig:regression}

\subsection{Final Test Results}

\begin{table}[H]
\centering
\begin{tabular}{@{}lrrr@{}}
\toprule
\textbf{Category} &
\textbf{Total} &
\textbf{Passed} &
\textbf{Success Rate}\\
\midrule
Model unit tests & 10 & 10 & 100\%\\
API tests & 11 & 11 & 100\%\\
Integration tests & 2 & 2 & 100\%\\
\midrule
\textbf{Total} & \textbf{23} & \textbf{23} & \textbf{100\%}\\
\bottomrule
\end{tabular}
\caption{Final local automated test results.}
\end{table}

% =========================================================
% 9. STATIC CODE ANALYSIS
% =========================================================

\section{Static Code Analysis}

Ruff was used to perform static code analysis and
identify Python code quality issues \cite{ruff}.

The initial analysis identified seven import-order
issues.

The supported automatic corrections were applied using:

\begin{lstlisting}
python -m ruff check api/ ml_model/ worker/ tests/ --fix
\end{lstlisting}

A subsequent analysis was performed using:

\begin{lstlisting}
python -m ruff check api/ ml_model/ worker/ tests/
\end{lstlisting}

The final output was:

\begin{lstlisting}
All checks passed!
\end{lstlisting}

The complete test suite was rerun successfully
following the lint corrections.

\shot{30_ruff_linting.png}
{Successful Ruff static code analysis.}
{fig:ruff}

Ruff was also executed successfully after the latest
infrastructure configuration changes and local
environment migration.

% =========================================================
% 10. VERSION CONTROL
% =========================================================

\section{Version Control and GitHub}

\subsection{Repository Initialization}

Git was used to track source code, infrastructure
configuration, tests and documentation.

The repository was initialized locally and the
primary branch was named \texttt{main}.

\begin{lstlisting}
git init
git branch -M main
git add .
git commit -m "Initial commit: ML system"
git push -u origin main
\end{lstlisting}

The initial commit was:

\begin{lstlisting}
c6ae85a
\end{lstlisting}

\shot{31_git_initial_commit.png}
{Initial Git commit and repository history.}
{fig:git}

\subsection{GitHub Repository}

The project was published to the following public
GitHub repository:

\begin{center}
\url{https://github.com/neskrotna/sepl-ml-system}
\end{center}

\shot{32_github_repository.png}
{Published GitHub project repository.}
{fig:repo}

\subsection{Version Control Exclusions}

The \codepath{.gitignore} file excludes local development
artifacts, including virtual environments, Python
cache files and the actual environment configuration.

The configuration template remains version-controlled.

The exclusions were verified using:

\begin{lstlisting}
git check-ignore .env
git check-ignore .venv/
\end{lstlisting}

Both commands returned the corresponding ignored
file or directory.

This prevents machine-specific environment files
from being included in repository commits.

\subsection{Infrastructure Configuration Commit}

The updated Docker configuration, environment template,
GitHub Actions workflow, README and implementation
screenshots were committed to GitHub.

The commit message was:

\begin{lstlisting}
feat: add model bind mounts and environment configuration
\end{lstlisting}

The verified commit identifier was:

\begin{lstlisting}
a102b16
\end{lstlisting}

Pushing this commit to the main branch automatically
triggered the updated continuous integration workflow.

% =========================================================
% 11. CONTINUOUS INTEGRATION
% =========================================================

\section{Continuous Integration}

\subsection{GitHub Actions Workflow}

Continuous integration was implemented using
GitHub Actions \cite{actions}.

The workflow is defined in:

\begin{lstlisting}
.github/workflows/ci.yml
\end{lstlisting}

The workflow is named:

\begin{lstlisting}
White Wine Quality CI
\end{lstlisting}

It is triggered by:

\begin{itemize}
    \item Pushes to the main branch.
    \item Pull requests targeting the main branch.
    \item Manual workflow execution.
\end{itemize}

The initial workflow commit was \texttt{fbe8b79}.

\subsection{CI Job Structure}

The workflow consists of two independent jobs.

\subsubsection*{Job 1: Lint and Unit Tests}

The first job performs the following steps:

\begin{enumerate}
    \item Check out the repository.
    \item Set up Python 3.12.
    \item Install project dependencies.
    \item Execute Ruff static code analysis.
    \item Execute the 21 non-integration tests.
\end{enumerate}

\subsubsection*{Job 2: Docker Integration Tests}

The updated second job performs the following steps:

\begin{enumerate}
    \item Check out the repository.
    \item Create the environment configuration.
    \item Set up Python 3.12.
    \item Install test dependencies.
    \item Build and start Docker Compose services.
    \item Wait for API readiness.
    \item Execute the two integration tests.
    \item Display Docker logs if a failure occurs.
    \item Stop Docker Compose services.
\end{enumerate}

\subsection{CI Execution Architecture}

\begin{center}
\fbox{
\begin{minipage}{.9\linewidth}
\centering
\vspace{5pt}

\textbf{Push / Pull Request / Manual Trigger}\\[5pt]
$\downarrow$\\[5pt]
\textbf{GitHub Actions}\\[8pt]

\begin{tabularx}{\linewidth}{@{}XX@{}}
\centering\textbf{Job 1} &
\centering\arraybackslash\textbf{Job 2}\\

\centering Lint and Unit Tests &
\centering\arraybackslash Docker Integration Tests\\

\centering $\downarrow$ &
\centering\arraybackslash $\downarrow$\\

\centering Ruff &
\centering\arraybackslash Environment configuration\\

\centering $\downarrow$ &
\centering\arraybackslash $\downarrow$\\

\centering 21 non-integration tests &
\centering\arraybackslash Docker startup\\

&
\centering\arraybackslash $\downarrow$\\

&
\centering\arraybackslash API readiness check\\

&
\centering\arraybackslash $\downarrow$\\

&
\centering\arraybackslash 2 integration tests\\
\end{tabularx}

\vspace{6pt}

$\downarrow$\\[4pt]
\textbf{CI Workflow Result}

\vspace{5pt}
\end{minipage}
}
\end{center}

\subsection{Initial GitHub Account Restriction}

The first hosted workflow execution was prevented
from starting because of an account-level billing
restriction.

GitHub reported:

\begin{quote}
\textit{The job was not started because your account
is locked due to a billing issue.}
\end{quote}

Neither CI job started during the affected execution.

The restriction was therefore unrelated to application
code, automated test failures or Docker configuration.

Local automated tests and containerized integration
were verified independently.

\shot{33_github_actions_billing_block.png}
{Historical GitHub account billing restriction preventing CI execution.}
{fig:ci}

\subsection{Resolution of the Billing Restriction}

GitHub Support was contacted to investigate the
account restriction.

On 24 November 2026, GitHub Support confirmed that
an expired coupon had triggered a billing-related
account lock.

The account was unlocked and its billing state
was reset.

Following the resolution, the GitHub Actions
workflow was executed again.

\subsection{First Successful Hosted CI Execution}

Workflow run number 3 completed successfully on
24 November 2026.

The workflow was manually triggered on the main
branch.

The associated commit was \texttt{c450be8}.

\begin{table}[H]
\centering
\begin{tabularx}{\linewidth}{@{}Xll@{}}
\toprule
\textbf{CI Job} & \textbf{Duration} & \textbf{Result}\\
\midrule
Lint and Unit Tests & 37 seconds & Success\\
Docker Integration Tests & 1 min 35 sec & Success\\
\midrule
Total workflow duration & 1 min 39 sec & Success\\
\bottomrule
\end{tabularx}
\caption{First successful hosted CI workflow.}
\end{table}

The successful execution confirmed that the previous
account-level restriction had been resolved.

The workflow is available at:

\begin{center}
\url{https://github.com/neskrotna/sepl-ml-system/actions/runs/45367264475}
\end{center}

\optionalshot{34_github_actions_success.png}
{First successful GitHub Actions workflow execution.}
{fig:firstcisuccess}

\optionalshot{35_ci_lint_unit_tests.png}
{Successful linting and non-integration testing in GitHub Actions.}
{fig:firstciunit}

\optionalshot{36_ci_docker_integration_tests.png}
{Successful Docker integration testing in GitHub Actions.}
{fig:firstciintegration}

The workflow overview also displayed non-blocking
Node.js deprecation annotations associated with
GitHub Actions dependencies.

Both jobs completed successfully despite these
annotations.

\subsection{CI Environment Configuration Update}

After introducing the local environment file, the
Docker integration job was extended to generate its
own runtime configuration.

The following step was added immediately after
repository checkout:

\begin{lstlisting}
- name: Create environment configuration
  run: cp .env.example .env
\end{lstlisting}

This is required because the actual \codepath{.env}
file is excluded from version control.

The GitHub-hosted runner therefore creates its own
configuration from the committed template before
starting Docker Compose.

The isolated lint and unit testing job does not
require this additional step.

\subsection{API Readiness Verification}

The integration workflow waits for the API to
become available before executing integration tests.

The readiness check repeatedly requests:

\begin{lstlisting}
http://localhost:8000/health
\end{lstlisting}

The workflow allows up to 30 attempts with a
two-second delay between unsuccessful attempts.

If the API does not become ready, Docker logs
are displayed and the job fails.

This prevents integration tests from starting
before the application is available.

\subsection{Successful Updated CI Execution}

The updated infrastructure and CI workflow were
pushed to the main branch on 24 November 2026.

Workflow run number 5 was automatically triggered
by the push.

Both jobs completed successfully.

\begin{table}[H]
\centering
\begin{tabularx}{\linewidth}{@{}lX@{}}
\toprule
\textbf{Property} & \textbf{Observed Value}\\
\midrule
Workflow & White Wine Quality CI\\
Run number & 5\\
Branch & main\\
Commit & a102b16\\
Trigger & Push\\
Overall status & Success\\
Successful jobs & 2/2\\
Total duration & 1 minute 40 seconds\\
\bottomrule
\end{tabularx}
\caption{Updated GitHub Actions workflow execution.}
\end{table}

\begin{table}[H]
\centering
\begin{tabularx}{\linewidth}{@{}Xll@{}}
\toprule
\textbf{CI Job} & \textbf{Duration} & \textbf{Result}\\
\midrule
Lint and Unit Tests & 29 seconds & Success\\
Docker Integration Tests & 1 min 37 sec & Success\\
\bottomrule
\end{tabularx}
\caption{Successful execution of both updated CI jobs.}
\end{table}

\shot{43_github_actions_updated_configuration.png}
{Successful automatically triggered CI workflow after infrastructure changes.}
{fig:updatedci}

\optionalshot{44_ci_updated_unit_tests.png}
{Updated CI workflow linting and non-integration test results.}
{fig:updatedciunit}

\optionalshot{45_ci_updated_integration_tests.png}
{Updated CI workflow Docker integration test results.}
{fig:updatedciintegration}

\subsection{Continuous Integration Verification}

The successful automatically triggered workflow
demonstrates that the updated infrastructure works
on GitHub-hosted runners.

The workflow verifies application quality through
static analysis, isolated tests and containerized
integration testing.

The CI jobs run independently and their execution
durations therefore do not need to add up to the
total workflow duration.

The implementation provides continuous integration.
Automated deployment to a production environment
has not been implemented.

% =========================================================
% 12. LOGGING AND MONITORING
% =========================================================

\section{Application Logging and Monitoring}

\subsection{Logging Implementation}

The API and background worker use Python's built-in
logging module \cite{logging}.

Logging supports observation of application behaviour,
prediction processing and error diagnosis.

The application log format includes:

\begin{itemize}
    \item Timestamp.
    \item Logging level.
    \item Descriptive event message.
\end{itemize}

An example application log entry is:

\begin{lstlisting}
2026-11-24 10:14:31,793 | INFO | Starting processing for task ...
\end{lstlisting}

\subsection{API Logging}

The API records important application events,
including:

\begin{itemize}
    \item Application startup.
    \item Redis connection establishment.
    \item Prediction task submission.
    \item Prediction status retrieval.
    \item HTTP requests and response status codes.
    \item Application errors.
\end{itemize}

API logs were inspected using:

\begin{lstlisting}
docker compose logs api --tail 30
\end{lstlisting}

The observed logs confirmed successful application
startup, Redis connectivity, health-check requests,
prediction submission and result retrieval.

\shot{41b_api_logging.png}
{API logs showing startup, health checks, prediction submission and result retrieval.}
{fig:apilogging}

\subsection{Worker Logging}

The background worker records important processing
events, including:

\begin{itemize}
    \item Worker startup.
    \item Successful model loading.
    \item Redis connection establishment.
    \item Waiting for prediction jobs.
    \item Task reception.
    \item Processing start.
    \item Completed predictions.
    \item Prediction confidence.
    \item Inference execution time.
\end{itemize}

Worker logs were inspected using:

\begin{lstlisting}
docker compose logs worker --tail 30
\end{lstlisting}

\shot{41a_worker_logging.png}
{Worker logs showing model loading, task processing and inference execution time.}
{fig:workerlogging}

\subsection{Task Traceability}

The API and worker logs contain the same unique task
identifier for an individual prediction.

During verification, the following task identifier
appeared in both service logs:

\begin{lstlisting}
c6ef8af3-80a3-4dd3-ba30-eaafafc019ae
\end{lstlisting}

The API logged the creation of the prediction job.

The worker subsequently logged reception, processing
and completion of the same task.

The API then logged retrieval of the task result.

This supports tracing a prediction across the
distributed application services.

\subsection{Health Monitoring}

The API exposes the following endpoint:

\begin{lstlisting}
GET /health
\end{lstlisting}

The endpoint checks Redis connectivity and reports
the corresponding application health state.

Redis availability is additionally monitored using
the Docker Compose health check.

These mechanisms provide basic service availability
monitoring.

The API health endpoint does not independently verify
all worker functionality or model prediction quality.

\subsection{Live Application Monitoring}

Docker Compose provides access to application logs
without requiring direct access to the container
filesystem.

Combined logs are available using:

\begin{lstlisting}
docker compose logs
\end{lstlisting}

Continuous monitoring of the API and worker is
available using:

\begin{lstlisting}
docker compose logs -f api worker
\end{lstlisting}

This supports runtime troubleshooting during
application development and testing.

\subsection{Observed Inference Execution Time}

The worker logs record the execution time associated
with prediction processing.

During local verification on 24 November 2026,
one completed prediction produced the following
observed values:

\begin{table}[H]
\centering
\begin{tabularx}{\linewidth}{@{}lX@{}}
\toprule
\textbf{Metric} & \textbf{Observed Value}\\
\midrule
Prediction & \texttt{good\_quality}\\
Confidence & 0.7500\\
Inference execution time & 0.3005 seconds\\
Task status & Completed\\
\bottomrule
\end{tabularx}
\caption{Observed prediction processing measurement.}
\end{table}

The recorded inference time corresponds to
approximately 300.5 milliseconds.

This measurement demonstrates successful runtime
performance logging.

It represents one observed prediction and should
not be interpreted as an average latency, a
throughput benchmark or a complete end-to-end
response time measurement.

\subsection{Monitoring Limitations}

The implemented monitoring approach provides
application logs, task traceability, health checks
and inference execution-time measurements.

Dedicated metrics aggregation, dashboards,
centralized log storage and automated alerting
have not been implemented.

These mechanisms represent potential extensions
for a production-oriented deployment.

% =========================================================
% 13. REPRODUCIBILITY AND DOCUMENTATION
% =========================================================

\section{Reproducibility and Documentation}

\subsection{Application Startup}

The repository contains the trained model artifact,
dataset, application source code, dependency
definitions and infrastructure configuration.

The recommended application startup procedure
requires Git and Docker Compose.

\begin{lstlisting}
git clone https://github.com/neskrotna/sepl-ml-system.git

cd sepl-ml-system

Copy-Item .env.example .env

docker compose up -d --build

docker compose ps
\end{lstlisting}

The application is accessible through:

\begin{center}
\url{http://localhost:8000/docs}
\end{center}

Retraining the model is not required for normal
application inference.

\subsection{Application Shutdown}

The Docker application is stopped using:

\begin{lstlisting}
docker compose down
\end{lstlisting}

The named Redis volume remains available after
normal shutdown.

\subsection{README Documentation}

The repository includes a comprehensive
\codepath{README.md} describing:

\begin{itemize}
    \item Project overview and architecture.
    \item Dataset and model results.
    \item Required development tools.
    \item Environment configuration.
    \item Docker startup instructions.
    \item Model bind mount configuration.
    \item API prediction examples.
    \item Logging and monitoring commands.
    \item Local Python environment setup.
    \item Model training instructions.
    \item Automated testing commands.
    \item Ruff static code analysis.
    \item GitHub Actions workflow.
    \item Troubleshooting procedures.
    \item Project implementation status.
\end{itemize}

The README supports reproduction of the
implemented application on another development
computer.

\subsection{Project Structure}

\begin{lstlisting}
sepl-ml-system/
|-- .github/
|   `-- workflows/
|       `-- ci.yml
|-- api/
|   |-- main.py
|   `-- schemas.py
|-- data/
|   `-- winequality-white.csv
|-- docs/
|   |-- report/
|   `-- screenshots/
|-- ml_model/
|   |-- explore_data.py
|   |-- model.py
|   |-- model.pkl
|   |-- test_local_prediction.py
|   `-- train.py
|-- tests/
|   |-- test_api.py
|   |-- test_integration.py
|   `-- test_model.py
|-- worker/
|   `-- worker.py
|-- .dockerignore
|-- .env.example
|-- .gitignore
|-- Dockerfile
|-- compose.yaml
|-- pytest.ini
|-- requirements.txt
`-- README.md
\end{lstlisting}

The directory structure separates application
components, machine learning development, tests,
infrastructure and documentation.

% =========================================================
% 14. HPC SIMULATION WITH SLURM
% =========================================================

\section{HPC Simulation with Slurm}

\textbf{Implementation status: Outstanding.}

The HPC simulation represents a subsequent
assessment component.

The existing application provides containerized
asynchronous inference using Redis and a
background worker.

This architecture is distinct from an HPC
batch scheduling environment.

The Slurm implementation, execution procedure
and experimental results will be documented
after the corresponding assessment requirements
have been implemented and verified.

No Slurm execution or HPC benchmark results
are claimed at the current project stage.

% =========================================================
% 15. PERFORMANCE REFLECTION
% =========================================================

\section{Performance Reflection}

\textbf{Implementation status: Preliminary observations.}

\subsection{Current Performance Observations}

The project currently provides the following
verified performance-related measurements.

\begin{table}[H]
\centering
\begin{tabularx}{\linewidth}{@{}Xl@{}}
\toprule
\textbf{Measurement} & \textbf{Observed Value}\\
\midrule
Model test accuracy & 75.16\%\\
Macro-averaged F1-score & Approximately 0.72\\
Serialized model size & Approximately 7.47 MB\\
Single observed inference time & 0.3005 seconds\\
Isolated model-loading test execution & 8.60 seconds\\
Latest CI lint and unit testing job & 29 seconds\\
Latest CI integration testing job & 1 min 37 sec\\
Latest complete CI workflow & 1 min 40 sec\\
\bottomrule
\end{tabularx}
\caption{Currently available performance-related measurements.}
\end{table}

The model-loading test duration includes test
execution overhead and is not an isolated
measurement of model deserialization time.

Similarly, GitHub Actions execution durations
include infrastructure preparation, dependency
installation and testing. They do not represent
model inference latency.

\subsection{Asynchronous Architecture Reflection}

The asynchronous design separates request
acceptance from prediction execution.

This supports independent management of the
API and background worker.

However, Redis communication, job serialization,
queue management and result retrieval introduce
additional processing overhead.

For a relatively lightweight Random Forest model,
direct synchronous inference might provide lower
latency for individual requests.

The selected architecture primarily demonstrates
asynchronous processing and separation of
application responsibilities.

A controlled comparison between synchronous and
asynchronous inference has not yet been performed.

\subsection{Pending Performance Evaluation}

The final performance reflection will be extended
after completion of the HPC simulation and the
remaining assessment requirements.

The currently available single-prediction
measurement is insufficient for conclusions
about average latency, throughput, scalability
or HPC performance.

% =========================================================
% 16. DISCUSSION AND LIMITATIONS
% =========================================================

\section{Discussion and Limitations}

\subsection{Implementation Outcomes}

The project demonstrates an end-to-end machine
learning engineering workflow.

The trained Random Forest model achieved 75.16\%
accuracy on the held-out test dataset.

The model was successfully integrated into a
validated REST API and asynchronous background
processing architecture.

Docker Compose provides a reproducible
multi-service environment.

The updated infrastructure supports externalized
configuration, read-only model access, persistent
Redis storage and service startup coordination.

Automated testing and static analysis were
successfully executed locally and through
GitHub Actions.

The application also provides basic runtime
observability through logging, health checks,
task identifiers and inference-time measurements.

\subsection{Model Limitations}

The model predicts binary quality categories
rather than the original numerical wine quality
score.

Evaluation was performed using a single
training and test split.

Generalization to wine samples collected from
different data sources or production conditions
has not been independently evaluated.

The reported confidence values have not been
evaluated through a dedicated probability
calibration experiment.

\subsection{Infrastructure Limitations}

Redis provides a lightweight queue and result
store suitable for the selected assessment use
case.

Advanced retry policies, dead-letter queues
and stronger delivery guarantees have not been
implemented.

The application uses a single configured
background worker.

Horizontal scaling and concurrent workload
performance have not been systematically
evaluated.

\subsection{Monitoring Limitations}

The application implements basic logging and
health monitoring.

Dedicated monitoring dashboards, centralized
log aggregation, automated alerting and model
drift detection remain potential future
extensions.

The observed inference-time measurement
represents one prediction rather than a
statistically evaluated workload.

\subsection{Environment Reproducibility}

The Docker and CI environments use Python 3.12,
while the recreated local development environment
uses Python 3.13.14.

All local tests passed in the recreated
environment.

Nevertheless, using consistent Python and
dependency versions across training, local
development, Docker and CI would further
reduce compatibility risks.

Serialized scikit-learn models should be loaded
using compatible library versions.

\subsection{Deployment and Security Limitations}

The API is exposed locally through the
loopback interface.

Authentication, authorization, production TLS
configuration and automated production
deployment have not been implemented.

The GitHub Actions workflow provides continuous
integration rather than a complete production
continuous deployment pipeline.

\subsection{Future Work}

Potential extensions include:

\begin{itemize}
    \item Cross-validation and model comparison.
    \item Hyperparameter optimization.
    \item Probability calibration.
    \item Model versioning.
    \item Improved dependency version management.
    \item Advanced Redis retry mechanisms.
    \item Multiple inference workers.
    \item Centralized metrics and log collection.
    \item Automated alerting.
    \item Authentication and authorization.
    \item Automated production deployment.
    \item Controlled latency and throughput benchmarking.
\end{itemize}

The outstanding Slurm simulation and final
performance reflection will be addressed as
part of the remaining assessment work.

% =========================================================
% 17. ASSESSMENT IMPLEMENTATION STATUS
% =========================================================

\section{Assessment Implementation Status}

The following table summarizes the current
implementation and verification status.

\begin{table}[H]
\centering
\small
\begin{tabularx}{\linewidth}{@{}cXX@{}}
\toprule
\textbf{Part} &
\textbf{Assessment Component} &
\textbf{Current Status}\\
\midrule

1 &
ML Model Development &
Implemented and verified\\

2 &
FastAPI Inference API &
Implemented and verified\\

3 &
Background Worker &
Implemented and verified\\

4 &
Containerization &
Verified locally and in CI\\

5 &
Testing and Quality Assurance &
23 local tests passed; both CI jobs successful\\

6 &
Logging and Monitoring &
Basic logging, health checks and task traceability verified\\

7 &
HPC Simulation with Slurm &
Outstanding\\

8 &
Performance Reflection &
Preliminary observations documented\\

9 &
Final Deliverables &
Documentation in progress\\

\bottomrule
\end{tabularx}
\caption{Current assessment implementation status.}
\end{table}

The status reflects the implementation evidence
available on 24 November 2026.

The final assessment completion status remains
subject to the outstanding implementation work
and verification against the complete assessment
requirements.

% =========================================================
% 18. CONCLUSION
% =========================================================

\section{Conclusion}

A functioning white wine quality prediction system
was implemented using Random Forest, FastAPI,
Pydantic, Redis, a Python background worker and
Docker Compose.

The machine learning model achieved a test
accuracy of 75.16\% and was successfully
serialized for reusable application inference.

The REST API validates incoming requests and
submits asynchronous prediction jobs to Redis.
The background worker processes queued jobs and
stores completed predictions for retrieval.

The containerized infrastructure was extended
with environment-based configuration, read-only
model bind mounts, persistent Redis storage,
service networking and health checks.

Runtime verification confirmed successful
container startup, Redis connectivity, model
access and end-to-end prediction processing.

All 23 automated tests passed locally, and
Ruff reported no remaining lint issues.

The GitHub account restriction that initially
prevented hosted CI execution was resolved.

Both the manually triggered workflow and the
subsequent automatically triggered workflow
completed successfully on GitHub-hosted runners.

The latest verified CI workflow completed in
1 minute and 40 seconds, with both jobs passing.

Application logging provides visibility into
startup events, task processing, prediction
results and inference execution time.

The current implementation demonstrates the
integration of machine learning development
with reproducible infrastructure, asynchronous
application architecture, automated testing,
continuous integration and basic runtime
monitoring.

The remaining assessment work consists of
the HPC simulation with Slurm, the final
performance reflection and completion of the
submission documentation.

% =========================================================
% REFERENCES
% =========================================================

\clearpage

\addcontentsline{toc}{section}{References}

\begin{thebibliography}{12}

\bibitem{wine}
Cortez, P., Cerdeira, A., Almeida, F.,
Matos, T. and Reis, J. (2009).
Modeling wine preferences by data mining
from physicochemical properties.
\emph{Decision Support Systems},
47(4), 547--553.

\bibitem{uci}
UCI Machine Learning Repository.
Wine Quality Dataset.
\url{https://archive.ics.uci.edu/dataset/186/wine+quality}

\bibitem{sklearn}
scikit-learn Documentation.
\url{https://scikit-learn.org/stable/}

\bibitem{fastapi}
FastAPI Documentation.
\url{https://fastapi.tiangolo.com/}

\bibitem{pydantic}
Pydantic Documentation.
\url{https://docs.pydantic.dev/}

\bibitem{redis}
Redis Documentation.
\url{https://redis.io/docs/}

\bibitem{docker}
Docker Documentation.
\url{https://docs.docker.com/}

\bibitem{compose}
Docker Compose Documentation.
\url{https://docs.docker.com/compose/}

\bibitem{pytest}
pytest Documentation.
\url{https://docs.pytest.org/}

\bibitem{ruff}
Ruff Documentation.
\url{https://docs.astral.sh/ruff/}

\bibitem{actions}
GitHub Actions Documentation.
\url{https://docs.github.com/en/actions}

\bibitem{logging}
Python Logging Documentation.
\url{https://docs.python.org/3/library/logging.html}

\end{thebibliography}

\end{document}